\documentclass[11pt,a4paper]{article}

\usepackage[a4paper,top=1.65cm,bottom=1.55cm,left=1.75cm,right=1.75cm]{geometry}
\usepackage[T1]{fontenc}
\usepackage{helvet}
\renewcommand{\familydefault}{\sfdefault}
\usepackage{xcolor}
\usepackage{hyperref}
\usepackage{tabularx}
\usepackage{enumitem}
\usepackage{ragged2e}
\usepackage{tikz}

\definecolor{navy}{HTML}{17233B}
\definecolor{blue}{HTML}{4568F2}
\definecolor{muted}{HTML}{52617A}

\hypersetup{colorlinks=true,urlcolor=blue,linkcolor=blue}
\pagestyle{empty}
\setlength{\parindent}{0pt}
\setlength{\parskip}{0pt}
\linespread{1.08}
\setlist[itemize]{leftmargin=1.2em,nosep,topsep=4pt}

\newcommand{\cvsection}[1]{%
  \vspace{13pt}%
  {\large\bfseries\color{navy}#1}\par
  \vspace{3pt}{\color{blue}\rule{\linewidth}{0.8pt}}\vspace{5pt}
}

% Small vector icons; no icon-font package is required.
\newcommand{\iconframe}[1]{%
  \tikz[baseline=-0.55ex]{%
    \draw[blue,line width=.55pt,rounded corners=.025cm]
      (0,0) rectangle (.27,.27);
    #1
  }%
}
\newcommand{\iconlocation}{\iconframe{%
  \draw[blue,line width=.6pt]
    (.135,.045) .. controls (.04,.14) and (.075,.225) .. (.135,.225)
    .. controls (.195,.225) and (.23,.14) .. (.135,.045);
  \fill[blue] (.135,.16) circle (.018);
}}
\newcommand{\iconphone}{\iconframe{%
  \draw[blue,line width=.6pt,rounded corners=.01cm]
    (.085,.045) rectangle (.185,.225);
  \fill[blue] (.135,.065) circle (.007);
}}
\newcommand{\iconmail}{\iconframe{%
  \draw[blue,line width=.6pt] (.045,.075) rectangle (.225,.195);
  \draw[blue,line width=.6pt]
    (.045,.195)--(.135,.13)--(.225,.195);
}}
\newcommand{\iconweb}{\iconframe{%
  \draw[blue,line width=.6pt] (.135,.135) circle (.085);
  \draw[blue,line width=.6pt] (.05,.135)--(.22,.135);
  \draw[blue,line width=.6pt] (.135,.05) ellipse (.035 and .085);
}}
\newcommand{\iconcode}{\iconframe{%
  \draw[blue,line width=.7pt] (.1,.075)--(.055,.135)--(.1,.195);
  \draw[blue,line width=.7pt] (.17,.075)--(.215,.135)--(.17,.195);
}}
\newcommand{\iconnetwork}{\iconframe{%
  \node[text=blue,font=\fontsize{6}{6}\selectfont\bfseries]
    at (.135,.135) {in};
}}
\newcommand{\contact}[2]{#1\hspace{6pt}#2}

\begin{document}
\RaggedRight

{\fontsize{25}{28}\selectfont\bfseries\color{navy}Arjun Sharma}\par
\vspace{4pt}
{\normalsize\bfseries\color{blue}BSc Computer Science Student}\par
\vspace{9pt}

{\small\color{muted}
\begin{tabularx}{\linewidth}{@{}XXX@{}}
\contact{\iconlocation}{Berlin, Germany} &
\contact{\iconphone}{015901445590} &
\contact{\iconmail}{\href{mailto:Arjunhry2006@gmail.com}{Arjunhry2006@gmail.com}} \\[6pt]
\contact{\iconweb}{\href{https://arjunsh4198-source.github.io/arjun-portfolio-code/}{Portfolio}} &
\contact{\iconcode}{\href{https://github.com/arjunsh4198-source}{GitHub}} &
\contact{\iconnetwork}{\href{https://www.linkedin.com/in/arjun-sharma-3b1b25358/}{LinkedIn}} \\
\end{tabularx}
}

\cvsection{Professional Summary}
BSc Computer Science student at Gisma University of Applied Sciences,
expected to graduate in 2028. Building a foundation in web development,
Python programming and databases through practical projects. Interested
in creating clear, useful software and continuing to learn through
hands-on work.

\cvsection{Education}
\textbf{BSc Computer Science}\hfill
{\color{muted}Expected graduation: 2028}\par
Gisma University of Applied Sciences\par
{\color{muted}Currently in the second semester.}

\cvsection{Technical Skills}
\begin{tabularx}{\linewidth}{@{}>{\bfseries\raggedright\arraybackslash}p{4.0cm}X@{}}
Web development & HTML, CSS, JavaScript \\
Programming \& data & Python, pandas, matplotlib \\
Databases & SQL, MySQL \\
Development tools & Git, GitHub, GitHub Pages, VS Code, LaTeX \\
\end{tabularx}

\cvsection{Projects}
\textbf{Personal Computer Science Portfolio}\par
Built a responsive website to present education, skills and projects.
Used HTML, CSS and JavaScript for the layout, mobile navigation and
project filtering, alongside an animated SVG illustration.

\vspace{7pt}
\textbf{Expense Tracker}\hfill{\color{muted}Python}\par
Created a command-line tool that records expenses in CSV format,
summarises spending with pandas and saves a category chart using
matplotlib.

\vspace{7pt}
\textbf{Library Database}\hfill{\color{muted}MySQL}\par
Designed related tables for books, members and loans, with sample data
and SQL queries for viewing current loans and borrowing activity.

\cvsection{Interests}
Reading books \quad\textbullet\quad Cricket \quad\textbullet\quad Coding
\quad\textbullet\quad Chess \quad\textbullet\quad Esports
\quad\textbullet\quad Music

\end{document}